% Act 1: fence the idea -- what is old (the three mechanisms) and what is new (what sharing requires).
\begin{frame}[t]{\ev{\tagO}\surf{WHAT IS NEW}Caches, forks and sandboxes are old; coupling them to evidence and acceptance is new.}
\vspace{.2cm}
{\fontsize{14}{17}\selectfont\renewcommand{\arraystretch}{1.35}
\begin{tabularx}{\linewidth}{@{}L{2.4cm}L{4.0cm}Y@{}}
\toprule
 & \textcolor{Teal}{\textbf{Keeps}} & \textcolor{Teal}{\textbf{Examples}}\\
\midrule
\textbf{Cache} & run outputs & \textcolor{Muted}{ccache, Bazel, Incredibuild}\\
\textbf{Fork} & reached states & \textcolor{Muted}{SnowFlock (Xen), CRIU, Firecracker}\\
\textbf{Sandbox} & writes private & \textcolor{Muted}{containers, gVisor, microVMs}\\
\bottomrule
\end{tabularx}\par}
\vspace{.15cm}
\claim{New: count shared evidence once; publish checked bytes once.}
\sources{Not claimed today: a fork speedup (our agent runs were serial) or live-memory restore. SnowFlock: Lagar-Cavilla et al., EuroSys 2009. Firecracker: Agache et al., NSDI 2020.}
\qadetail{The table names established mechanisms; no claim of invention, comprehensive priority or measured parallel speedup is made. Forking states to search is also old: KLEE (Cadar, Dunbar and Engler, OSDI 2008) and S2E (Chipounov, Kuznetsov and Candea, ASPLOS 2011) fork execution states; the AFL forkserver forks an initialized target per input; Potemkin (Vrable et al., SOSP 2005) flash-clones VMs. The claim is the evidence and acceptance contract around reuse, not search by forking. The contribution claim is the integrated execution, evidence and acceptance contract: reuse a state only when it beats rebuilding it, count shared observations once, and publish the exact checked bytes once from a controller. What each mechanism must guarantee: a cache key covers every input (only for hermetic steps: time, randomness, network and undeclared files must be declared or excluded); a fork's capture keeps what the next step needs; a sandbox hides undeclared inputs and keeps writes private. Fork semantics are backend-specific: worktrees preserve source views, process checkpoints (CRIU, gVisor save/restore) and VM snapshots (Xen, SnowFlock, Firecracker) preserve only their declared state surfaces, and filesystem snapshots do not establish memory capture. SnowFlock: Lagar-Cavilla et al., Rapid Virtual Machine Cloning for Cloud Computing, EuroSys 2009, built on Xen; Xen came from the Cambridge Computer Laboratory (Barham et al., SOSP 2003). Bazel here means its remote cache. A sandbox bounds local effects; it does not undo remote effects, and a restored child must not hold publication authority. Not claimed: our factory runs (SELFHOST-2 and SELFHOST-3, 27 September 2026) ran serially in Docker Sandboxes, not islo, and nothing forked, so no fork speedup is measured; we do not claim that a fork is faster than a good build cache. Restore fan-out means many restores of one named snapshot: a named 141 MB islo snapshot, restore--run--capture round trips, 255 of 256 succeeded at concurrency 12, p50 6.87 s, p95 9.04 s, client-measured, teardown excluded (arXiv:2607.09689v4, Table 2). These are API round trips, not mechanism latency. Memory preservation is unverified until the preregistered Gate 0 RAM-nonce probe runs, so no islo snapshot is described as a memory fork.}
\note{[0:50] A cache keeps run outputs, as in ccache, Bazel's remote cache or Incredibuild. A fork keeps reached states: SnowFlock cloned Xen machines in two thousand nine, and Xen came from this lab. A sandbox keeps writes private, as containers, gVisor and microVMs do. What is new is what sharing requires: count shared evidence once, and publish checked bytes once. I do not claim a fork speedup, since our agent runs were serial, or that our snapshots keep live memory. So what does sharing work cost us? Three rules.}
\end{frame}
